\begin{abstract}
Recent studies report that retrieving raw conversation history by embedding similarity matches dedicated memory systems on conversational benchmarks. We confirm this observation on text benchmarks once the retrieval budget is large, and we locate where it breaks in multimodal agent memory: many questions refer to a set of past items, and top-$k$ similarity ranking rarely returns the whole set. On MemEye, the share of questions whose complete evidence set appears among the top 10 rounds falls from 0.92 for single-round questions to 0.32 for sets of four or five rounds. We study query-time set selection (\qss{}), in which the memory keeps raw rounds, writes one short index line per round, and a language model reads the whole index at question time and returns the evidence set with a size that it decides. On MemEye, where we developed the method, \qss{} returns the complete set for 90.3\% of questions while opening 7.3 rounds, against 53.9\% for cosine retrieval of the same size, and same-model controls attribute most of this recall to model judgment and a recall-oriented instruction, while returning a set of model-chosen size adds a small further gain and lets the budget follow the size of the set. A pre-registered test on the held-out MemLens benchmark bounds these results: \qss{} still doubles the full-set recall of equal-size cosine retrieval, but its answers do not exceed those of top-40 retrieval. Two failures explain this gap. The index lines expose only 37\% of the question words that occur in MemLens evidence rounds, against 95\% on MemEye, and lightweight readers fail to aggregate evidence sets even when the sets are complete. Visibility-gated set reading (VGS) therefore estimates this index visibility at question time and, when it is low, widens the index lines and passes the reader per-member notes; in registered tests on held-out MemLens questions, it improves a lightweight reader by $+0.071$ over top-40 retrieval, and by $+0.056$ on histories twice as long, but on the external WorldMemArena benchmark, whose visibility is equally low, the frozen design gains no significant accuracy, so low visibility alone does not identify where the notes help. To control cost, a compact index halves the input of the selection call without lowering accuracy, and an explicit context cache stores this index once per user; with both, a strong reader with set selection is more accurate than top-40 retrieval on MemEye (0.803 against 0.749) at 64\% of its cost, an advantage that requires each cached index to serve enough questions. Query-time set selection is thus a conditional tool, and its conditions, namely index visibility, reader strength, and the persistence of the memory, can be measured and acted upon at question time.
\end{abstract}
